\section{Results}
\label{sec:experiments}

The 290 configurations leave three kinds of record, and this section reads each in turn: the scored
artifacts say what the agents \emph{achieved} (\S\ref{sec:results}), the submitted patches and audited
trajectories say what they \emph{did} (\S\ref{sec:behavior}), and the apparatus itself leaves a record
of how far its own conventions move the number (\S\ref{sec:analysis}) --- far enough that no result in
the first part can be read without the third.

Throughout, a configuration's score is the best of its up to three retained checkpoints under the
task's own metric direction: a harness convention fixed before any run, though what goes into the
three slots is the agent's to choose, and one worth $7.8$ points of pooled win rate on its own
(\S\ref{sec:ckpt}). ``Beat'' means strictly better than the reference, and we never average a raw
score across tasks because the ten metrics live in different units.

\subsection{What the agents achieve}
\label{sec:results}

We read the 792 scored artifacts three ways: two \emph{reference-free} aggregates that give the
standings among systems, a lifecycle reading of how many configurations reach a scored artifact at
all, and a reading against the recipe each repository ships, which is the one that says whether the
work is worth adopting.

\subsubsection{Standings: two reference-free aggregates}
\label{sec:standings}

Because the ten metrics are incommensurable and each reference is contestable, we rank systems by two
aggregates that read \emph{no} reference at all. On each task, the three best distinct configuration
scores observed across the six systems set its gold, silver and bronze thresholds; a system earns
$1.00/0.67/0.33$ for its best configuration reaching one, and its \emph{medal score} is the mean of
that credit over the ten tasks. Its \emph{average rank} is the mean of its within-task competition
rank, with a system that scores nothing on a task placed last. Neither construction divides by a start, a recipe, or a published
number, so both are defined identically on all ten tasks --- which is exactly why they, and not the
reference-relative rates of \S\ref{sec:pertask}, carry the standings.

Ordering ten tasks this way requires one scalar per artifact, and unlearning does not publish one:
its evaluation reports two components, and its task package forbids collapsing them into a proxy that
replaces them. The same package, however, supplies an internal scalar and designates it explicitly
\emph{for lifecycle ordering}. We follow that division. The aggregates here order artifacts, so they
use that scalar; the reference-relative rate of \S\ref{sec:pertask} reports an outcome, so it uses the
two components under dominance and never combines them (Appendix~\ref{app:unlearning}). Where the two
readings disagree they are answering different questions, and we report both rather than picking one.

Two of the ten tasks carry a qualification worth stating before any aggregate is read. On diffusion
fine-tuning the reward predictor keeps rising past the range its training labels span, so a higher
score there need not be a better image; on agentic RL the top scores come from imitating an optimal
solver the repository itself ships, so a higher score there separates recipe classes rather than
degrees of training. Both metrics are the repositories' own, and \S\ref{sec:validity} gives the
evidence. Every aggregate below is therefore reported twice: over all ten tasks, and over the eight
that remain when these two are removed. Code RL shows what a reference-free aggregate does and does not say. Its
gold threshold is $0.1328$ --- the best configuration in the field, which lands on exactly the shipped
recipe's own score of $17/128$ --- so the field tops out level with the repository and never above it.
Podium credit there records which systems reached that ceiling; the reference-relative rate of
\S\ref{sec:pertask} records that none cleared it, $0/28$, since ``beat'' means strictly better. Both
readings are correct because they answer different questions, and this is the task where the
difference is starkest.

\textbf{Claude Opus~5 leads both aggregates, though not by a separable margin}
(Table~\ref{tab:medal}): a medal score of $0.833$ and an average rank of $1.40$, a medal on nine of the
ten tasks and the top score on eight. Comparing Opus~5 against GPT-5.6 Sol task by task on each
system's best score gives five tasks to Opus, two to Sol and three ties --- a split a sign test cannot
distinguish from chance (one-sided $p=0.23$), on a design with no replication to set a noise floor
against. GPT-5.6 Sol is second on half that credit ($0.500$, rank $2.30$) and GPT-5.6 Terra third
($0.268$, $3.50$); Claude Sonnet~5 and GPT-5.6 Luna are close below them and exchange fourth and fifth
place depending on which aggregate is read ($0.167$/$0.133$, ranks $4.30$/$4.20$), and Kimi~K3 reaches
no medal. The eight-task columns leave the ranking intact, though two of Opus~5's eight golds sit on
the two qualified tasks.

\textbf{The ordering below first place is partly a statement about grid depth.} The three grids expose
different numbers of reasoning-effort levels --- six on Codex, five on Claude Code, one on Kimi~K3 ---
so a system with more levels gets more attempts at each podium. Restricting the pool to the single
\texttt{max} level every system ran leaves Claude Opus~5 first (medal $0.633$) but rearranges the field
beneath it: Claude Sonnet~5 rises to $0.401$ and Kimi~K3, which takes no medal at all in the full pool,
reaches $0.333$ --- ahead of GPT-5.6 Terra and Luna, though still behind Sol at $0.433$. Dropping only Codex's extra \texttt{none}
level, by contrast, leaves every medal score unchanged, so the effect is Kimi's single cell rather than
Codex's deeper grid. We report the full pool because it is the comparison the suite was run to
make, and this restriction because the full pool cannot separate a system from the number of attempts
it was given.

\begin{figure*}[t]
\centering
\includegraphics[width=0.88\linewidth]{fig_models}
\caption{\textbf{Model capability at a glance.} Each cell gives a system's rank ($1$--$6$) on a task,
shaded by the medal credit that rank actually earned
rather than by the rank itself: thresholds are set by the three best distinct \emph{configuration}
scores, which one system may hold several of, so a rank of three need not be a bronze and an unshaded
cell is a rank without a podium. Medal score is the bar at right, and rows are ordered by it. Claude Opus~5 ranks first on eight of the ten tasks and takes a podium on
nine; the four systems below it are separated mostly by single tasks.}
\label{fig:models}
\end{figure*}

\begin{table*}[t]
\centering
\caption{\textbf{The standings, by two reference-free aggregates.} On each task the three best
\emph{distinct} configuration scores across the six systems set its gold\,/\,silver\,/\,bronze
thresholds, and a system is credited $1.00/0.67/0.33$ for its best configuration reaching one. No
reference enters the construction, so both aggregates are defined on all ten tasks; the \textbf{8\,tk}
columns drop the two tasks \S\ref{sec:validity} qualifies, and the \textbf{\texttt{max}} columns
recompute the whole construction --- thresholds included --- on the single reasoning-effort level all
six systems ran, which is the only pool in which the grids are matched. \textbf{Claude Opus~5 leads every reading},
taking a medal on nine of ten tasks and ranking first on eight. Below it the ordering is not stable:
in the matched \texttt{max} pool Claude Sonnet~5 and Kimi~K3 rise (italic), the latter from no medal
at all to $0.333$, because the full pool gives Kimi one attempt per task against the others' five or
six. This podium is set by the six systems measured here and does not read any external
reference. Unlearning's thresholds are set on the internal scalar its task package designates for
lifecycle ordering; its reference-relative rate is scored by dominance on the two official components
instead (\S\ref{sec:standings}).}
\label{tab:medal}
\small
\begin{adjustbox}{max width=\linewidth}
\begin{tabular}{lcccccccccc rrr rrr}
\toprule
& \multicolumn{10}{c}{\textbf{medal credit per task}}
& \multicolumn{3}{c}{\textbf{medal} $\uparrow$}
& \multicolumn{3}{c}{\textbf{avg.\ rank} $\downarrow$} \\
\cmidrule(lr){2-11}\cmidrule(lr){12-14}\cmidrule(lr){15-17}
\textbf{System} & DDPO & DiGress & DPO & Soup & OPD & OpenR1 & NPO & OWL & RAGEN & BTRM
& 10\,tk & 8\,tk & \texttt{max} & 10\,tk & 8\,tk & \texttt{max} \\
\midrule
Claude Opus 5 & 1.00 & 0.33 & 1.00 & 0.00 & 1.00 & 1.00 & 1.00 & 1.00 & 1.00 & 1.00 & \textbf{0.833} & \textbf{0.791} & \textbf{0.633} & \textbf{1.40} & \textbf{1.50} & \textbf{2.60} \\
GPT-5.6 Sol & 0.00 & 1.00 & 0.00 & 1.00 & 0.00 & 1.00 & 1.00 & 0.00 & 1.00 & 0.00 & 0.500 & 0.500 & 0.433 & 2.30 & 2.38 & 3.10 \\
GPT-5.6 Terra & 0.00 & 0.67 & 0.00 & 0.00 & 0.67 & 0.67 & 0.00 & 0.00 & 0.67 & 0.00 & 0.268 & 0.251 & 0.267 & 3.50 & 3.50 & 3.90 \\
Claude Sonnet 5 & 0.00 & 0.00 & 0.00 & 0.67 & 0.00 & 0.67 & 0.00 & 0.00 & 0.33 & 0.00 & 0.167 & 0.168 & \emph{0.401} & 4.30 & 4.12 & \emph{3.10} \\
GPT-5.6 Luna & 0.00 & 0.00 & 0.33 & 0.00 & 0.00 & 0.33 & 0.00 & 0.00 & 0.00 & 0.67 & 0.133 & 0.166 & 0.300 & 4.20 & 3.88 & 4.00 \\
Kimi K3 & 0.00 & 0.00 & 0.00 & 0.00 & 0.00 & 0.00 & 0.00 & 0.00 & 0.00 & 0.00 & 0.000 & 0.000 & \emph{0.333} & 4.70 & 5.00 & \emph{3.70} \\
\midrule
\multicolumn{17}{l}{\emph{podium thresholds, in each task's own units and direction}} \\
\footnotesize gold & \footnotesize 17.67 & \footnotesize 63.72 & \footnotesize .6223 & \footnotesize .7006 & \footnotesize .4490 & \footnotesize .1328 & \footnotesize 1.00 & \footnotesize 12.96 & \footnotesize 1.00 & \footnotesize 77.09 & & & & & & \\
\footnotesize silver & \footnotesize 14.45 & \footnotesize 64.33 & \footnotesize .6150 & \footnotesize .6985 & \footnotesize .4427 & \footnotesize .1250 & \footnotesize .9960 & \footnotesize 13.04 & \footnotesize .9980 & \footnotesize 75.84 & & & & & & \\
\footnotesize bronze & \footnotesize 12.13 & \footnotesize 64.85 & \footnotesize .5642 & \footnotesize .6983 & \footnotesize .4401 & \footnotesize .1172 & \footnotesize .9929 & \footnotesize 13.25 & \footnotesize .8945 & \footnotesize 75.75 & & & & & & \\
\bottomrule
\end{tabular}
\end{adjustbox}
\par\vspace{2pt}
{\footnotesize Average rank uses competition ranking within a task; a system with no scored
configuration takes last place ($6$). Thresholds are rounded here and released at full evaluator
precision. Three restrictions bound what these columns may be compared to --- the podium is pool-defined; per-task credit is not constant, paying
$1.00$ where one system holds all three thresholds and up to $3.67$ where systems tie at them, so the
ten tasks do not weigh equally in the mean; and attempts per cell are unequal --- Kimi~K3 contributes one configuration per task where the others
contribute five or six, and at the single \texttt{max} level all six systems ran it reaches a medal
score of $0.333$ rather than $0.000$ (\S\ref{sec:standings}).}
\end{table*}


\subsubsection{Lifecycle survival}
\label{sec:lifecycle}

\begin{figure}[t]
\centering
\includegraphics[width=0.75\linewidth]{fig_lifecycle}
\caption{\textbf{Outcomes by system, each normalized to $100\%$.} Every bar is one system's full set
of configurations ($n$ at right), split into those that beat the shipped recipe (dark), produced a
scored artifact that did not beat it (light), or produced no scored artifact at all (warm grey).
Normalizing lets the mix be compared across systems of unequal size --- Kimi~K3's ten runs included.
Between adjacent systems the mix differs by less than the checkpoint
convention alone is worth (\S\ref{sec:analysis}), though across the whole field it spans $50\%$ to
$72\%$; no-artifact failures stay in the single digits to
low teens and fall almost entirely on the Codex grid ($18$ of the $19$).
Absolute survival is high --- 271 of 290 overall, 792 scored artifacts (\S\ref{sec:lifecycle}).}
\label{fig:lifecycle}
\end{figure}

Of the 290 configurations, 282 reach a formal terminal state and 271 produce at least one scored
artifact, for 792 in total (Figure~\ref{fig:lifecycle}).

\paragraph{What fails, and why.}
\label{sec:failures}
Nineteen configurations do not reach a scored artifact, and the failure modes are different in kind.
Eight are rejected at source extraction: our extractor could not recover a usable source-only patch,
so no formal run began. These are all on the Codex grid and seven of the eight sit at its two lowest
effort levels --- \emph{a low reasoning budget mostly produces work the harness cannot carry across
the visibility boundary, not work that trains and fails}, which is why the completion climb of
\S\ref{sec:effort} is a property of our tooling rather than of capability. The other eleven retrain
and terminate but leave no scorable checkpoint. A third failure is visible only inside a task: on
agentic RL one configuration scores exactly $0.000$ on all 512 evaluation boards at every one of its
three retained checkpoints, and two further configurations collapse to $0.000$ at their final
checkpoint. The zero-scoring configuration's logs show a response-format collapse --- the policy
emits an opening reasoning tag and then stops, so no action decodes and the board never advances ---
induced by a recipe that rewrote the prompt layer. That zero is a correct measurement of a policy
which has stopped producing parseable actions, not a training dynamic, and it is the reason a raw
mean over a task's configurations can be badly misleading.

\paragraph{Survival is high and even across systems.} Every Claude Opus~5 and every Kimi~K3 configuration produces an
artifact, and the Codex grid ranges from $87\%$ to $97\%$ (Table~\ref{tab:systems});
Appendix~\ref{app:lifegrid} resolves the outcome mix to each \emph{(model, effort)} cell. Because each
artifact is retrained from a pinned commit and scored by an evaluator the agent never writes to, a
surviving configuration is a result another party can reproduce, not a number the agent reports about
itself.

\begin{table}[t]
\centering
\caption{\textbf{The six evaluated systems and their lifecycle survival.} The unit under test is a
\emph{(model, native harness, effort)} triple; a model and its harness are not separable, so the
Codex and Claude~Code columns are not a clean model-only comparison. \textbf{Cfg} is the number of
configurations (tasks $\times$ effort levels); \textbf{Scored} is those producing an
independently scored artifact. \textbf{vs.\ recipe} counts, direction-aware, the configurations that
beat the repository's shipped recipe. Every task publishes one, so the column runs over all $271$
scored configurations; counts are printed rather than percentages so the denominator is never implied.
We report no pooled rate against the frozen starting checkpoint: two tasks disclaim that reference and
elsewhere it is an undertrained artifact (\S\ref{sec:results}). \textbf{Cost}
is total exploration API cost at each provider's official USD rates, rounded to the dollar per row so
that the rows sum to one dollar below the exact total, including Kimi~K3 at Moonshot's
\$3\,/\,\$15 per million input\,/\,output tokens (\$0.30 cached). \textbf{med.\ \$/cfg} is the median
per-configuration exploration cost, and is the comparable figure because the grids differ in depth: it
ranges roughly ten-fold across the six systems while win rates vary little, so exploration cost is
weakly coupled to per-task success (\S\ref{sec:cost}). The reference-free standings are in Table~\ref{tab:medal}.}
\label{tab:systems}
\small
\begin{adjustbox}{max width=\linewidth}
\begin{tabular}{llrrrrrr}
\toprule
\textbf{System} & \textbf{Harness} & \textbf{Efforts} & \textbf{Cfg} & \textbf{Scored}
& \textbf{vs.\ recipe} & \textbf{Cost} & \textbf{med.\ \$/cfg} \\
\midrule
Claude Opus 5   & Claude Code & 5 & 50 & 50 ($100\%$) & $36/50$ & \$909 & \$18 \\
Claude Sonnet 5 & Claude Code & 5 & 50 & 49 ($98\%$)  & $29/49$ & \$330 & \$6 \\
Kimi K3         & Claude Code & 1 & 10 & 10 ($100\%$) & $7/10$ & \$59 & \$6 \\
\midrule
GPT-5.6 Sol     & Codex & 6 & 60 & 58 ($97\%$) & $35/58$ & \$2{,}692 & \$45 \\
GPT-5.6 Terra   & Codex & 6 & 60 & 52 ($87\%$) & $30/52$ & \$870 & \$8 \\
GPT-5.6 Luna    & Codex & 6 & 60 & 52 ($87\%$) & $33/52$ & \$337 & \$4 \\
\midrule
\textbf{All}    &       &   & \textbf{290} & \textbf{271 ($93\%$)} & \textbf{$170/271$} & \$5{,}198 & \na{} \\
\bottomrule
\end{tabular}
\end{adjustbox}
\end{table}


\subsubsection{Reading each task against its own references}
\label{sec:pertask}

The medal and rank aggregates above cover all ten tasks and read no reference. A different question ---
is a given artifact worth adopting? --- is asked against the reference the task itself supplies: the repository's
\emph{shipped recipe} (\S\ref{sec:setup}). Every task has one, so this is the reading that covers the
whole scored corpus; because the ten metrics live in different units we pool rates rather than average
scores, and we report the per-task rates beside the pooled one throughout
(Figure~\ref{fig:capability}).

\begin{figure}[t]
\centering
\includegraphics[width=0.78\linewidth]{fig_capability}
\caption{\textbf{Per-task win rates against both recipe references.} Share of a task's successful
configurations that beat the shipped recipe $\rshort$ (orange) and the budget-matched twelve-hour
replay $\rlong$ (green, where it exists), under the task's own metric direction. Unlearning is scored
by dominance on its two official components rather than on a scalar (Appendix~\ref{app:unlearning}). Read on best checkpoints throughout, the replay is the \emph{harder} reference on three of the six
tasks that carry both (DDPO, DiGress, OPD) and the easier one on two (DPO, OpenR1), so the pooled rate barely moves --- a net seven configurations out of $164$,
though $69$ change verdict in one direction or the other --- while which measurement of the recipe one uses changes
the verdict in both directions (\S\ref{sec:validity}). Beating the shipped recipe is selective: the
per-task rates run from $0/28$ on code RL to $28/28$ on weight averaging, and the pooled figure is an
average over that spread rather than a description of any task in it.}
\label{fig:capability}
\end{figure}

We score improvement against the reference the question is actually about: the recipe each
repository itself ships. The frozen starting checkpoint is listed per task in Table~\ref{tab:tasks} but
carries no pooled rate. Two tasks disclaim it outright --- one-shot pruning requires removing $70\%$ of
the weights it is handed, so no compliant configuration can match its own start, and graph generation's
start is recorded by its own task package as an orientation figure measured under a different container
image --- and elsewhere the start is an undertrained artifact that twelve hours of retraining clears
almost by construction. A rate against it would measure the gap each task opened, not the agent.

\textbf{Against the shipped recipe, the picture splits by task.} Every task publishes a recipe, so
this comparison runs over the whole scored corpus: $170$ of $271$ configurations ($62.7\%$) beat it.
By task the median configuration wins on seven and loses on three, and all three losses are
tasks whose committed recipe is itself a strong, separately reported configuration: preference
optimization ($7/29$), code RL ($0/28$) and reward modelling ($9/29$). Two of the wins are softer than the pooled figure makes them appear: weight
averaging ships as its recipe the naive uniform average of all $72$ ingredients, and one-shot pruning
ships a reference nearly five times worse in perplexity than the dense artifact it prunes. Together
those two supply $55$ of the $170$ wins from $56$ configurations; removing them costs $9.2$ points. Two further tasks, diffusion fine-tuning and agentic RL, are the ones whose metric
\S\ref{sec:validity} flags; removing those instead costs $5.0$ points. A third softness is not visible
in any count: on on-policy distillation, where the rate reads $14/25$, not one of those wins
exceeds twice the evaluator's standard error, and on weight averaging $1$ of $28$ does
(\S\ref{sec:noise}). The reading the artifacts support is therefore narrower than ``agents improve AI systems'':
under a fixed twelve-hour budget, agents reliably recover a competent default and exceed a strong one
only on a minority of tasks --- the same story the reference-free standings tell.
